%! AP Statistics — Unit 9, Lesson 9.1 — Guided Notes (Answer Key)
\documentclass[10pt]{article}
\usepackage{apstats-article}
\usepackage{apstats-key}   % pulls in -boxes; do NOT also load -boxes

\begin{document}

\pageheader{Unit 9, Lesson 9.1}{Guided Notes --- Answer Key}
\namedateperiod

\begin{objectivebox}
By the end of this lesson, I will be able to:
\begin{itemize}
    \item[$\square$] \ansline{Explain why the sample slope $b$ is a statistic that varies by sample.}
    \item[$\square$] \ansline{State $H_0\colon \beta = 0$ and explain what it means in context.}
    \item[$\square$] \ansline{\smash{Ask whether an observed slope is convincing evidence that $\beta \ne 0$.}}
\end{itemize}
\end{objectivebox}

\begin{vocabbox}
\noindent\textbf{\textcolor{navy}{Population slope ($\beta$):}}\\[1pt]%
\ansline{The true rate of change of $y$ with $x$ in the population;}\\[1pt]\ansline{a fixed, unknown parameter.}\\[3pt]
\noindent\textbf{\textcolor{navy}{Sample slope ($b$):}}\\[1pt]%
\ansline{The slope of the LSRL computed from sample data;}\\[1pt]\ansline{a statistic that varies from sample to sample.}\\[3pt]
\noindent\textbf{\textcolor{navy}{Sampling variability:}}\\[1pt]%
\ansline{The natural variation in a statistic (such as $b$) across different}\\[1pt]\ansline{random samples from the same population.}\\[3pt]
\noindent\textbf{\textcolor{navy}{Null hypothesis for slope:}}\\[1pt]%
\ansline{$H_0\colon \beta = 0$ --- no linear relationship between $x$ and $y$ in the}\\[1pt]\ansline{\smash{population; any non-zero $b$ is due to chance.}}\\[3pt]
\end{vocabbox}

\begin{hookbox}
\textbf{Discussion:} Two scatter plots are shown on the board. Both have a
positive LSRL with slope $b > 0$. One has tightly clustered points; the other
has widely scattered points. Which trend are you more confident is real? Why?

\ansline{The tightly clustered one: with little scatter, a slope that big}\\\ansline{\smash{is unlikely to come from chance alone.}}
\end{hookbox}

\begin{notesbox}{The Sample Slope Is a Statistic}
In Units 2--3 we fit an LSRL $\hat{y} = a + bx$ to sample data and computed
the slope $b$. The slope $b$ is a \makebox[4cm][l]{\ans{statistic}}, not a parameter.

\vspace{0.1in}
\textbf{Key idea:} Different random samples from the same population produce
\makebox[5cm][l]{\ans{different}} values of $b$.

\vspace{0.1in}
The true (unknown) rate of change in the population is called the
\textbf{population slope}, denoted \makebox[1.5cm][l]{\ans{$\beta$}}.

\vspace{0.1in}
\renewcommand{\arraystretch}{1.6}
\begin{tabularx}{\linewidth}{Xcc}
    & \textbf{Symbol} & \textbf{Parameter or Statistic?} \\
    \hline
    True population slope       & $\beta$ & \makebox[4cm][l]{\ans{parameter}} \\
    Estimated slope from sample & $b$     & \makebox[4cm][l]{\ans{statistic}} \\
\end{tabularx}
\end{notesbox}

\begin{notesbox}{The Central Question of Unit 9}
When we see a non-zero slope $b$ in sample data, two explanations are possible:

\begin{enumerate}
  \item The true population slope $\beta$ \makebox[4cm][l]{\ans{$\ne 0$}} --- the trend is real.
  \item The true population slope $\beta = $ \makebox[1cm][l]{\ans{0}} --- there is no real
        linear relationship, and the non-zero $b$ is due to \makebox[4cm][l]{\ans{chance}}.
\end{enumerate}

\vspace{0.1in}
\textbf{The null hypothesis for slope:} $H_0\colon \beta = $ \makebox[1cm][l]{\ans{0}}

\vspace{0.08in}
This says: in the population, there is \makebox[5cm][l]{\ans{no linear relationship}}
between $x$ and $y$. Any non-zero $b$ we observe is just \makebox[4cm][l]{\ans{chance variation}}.

\vspace{0.1in}
\textbf{Unit 9 asks:} Is the sample slope $b$ we observed \makebox[5cm][l]{\ans{convincing}}
evidence that $\beta \ne 0$ in the population?
\end{notesbox}

\begin{notesbox}{Why Scatter Matters}
Two factors influence how surprising a non-zero $b$ is under $H_0\colon \beta = 0$:

\begin{enumerate}
  \item \textbf{Sample size ($n$):} A slope of $b = 3.4$ from $n = 200$ observations
        is \makebox[3cm][l]{\ans{stronger}} evidence than the same slope from $n = 10$ observations.
  \item \textbf{Scatter around the line:} The more tightly points cluster around the
        LSRL, the \makebox[3cm][l]{\ans{more precise}} our estimate of $\beta$.
\end{enumerate}

\vspace{0.1in}
This is why we can't just look at $b$ and declare the relationship real --- we need
\makebox[6cm][l]{\ans{a significance test}} to quantify how likely $b$ is under $H_0$.
\end{notesbox}

\begin{notesbox}{Unit 9 Arc Overview}
\renewcommand{\arraystretch}{1.5}
\begin{tabularx}{\linewidth}{lX}
    \textbf{Lesson 9.1} & Framing: what question does Unit 9 ask? \\
    \textbf{Lesson 9.2} & Sampling distribution of $b$; conditions for inference \\
    \textbf{Lesson 9.3} & Confidence interval for \makebox[2cm][l]{\ans{$\beta$}} \\
    \textbf{Lesson 9.4} & Significance test for \makebox[2cm][l]{\ans{$\beta$}} ($t$-test for slope) \\
    \textbf{Lesson 9.5} & Reading \makebox[4cm][l]{\ans{computer}} output for regression \\
    \textbf{Lesson 9.6} & Synthesis: \makebox[5cm][l]{\ans{choosing a procedure}} \\
\end{tabularx}
\end{notesbox}

\begin{practicebox}
\textbf{Guided Practice:} A study collects data on daily high temperature ($x$, in $^\circ$F)
and the number of ice cream cones sold ($y$) at a beachside stand for 20 days. The LSRL is
$\hat{y} = -82.4 + 3.4x$ with $r = 0.87$.

\begin{enumerate}
  \item Interpret the slope in context.
        \ansline{Each additional $1^\circ$F,}\\\ansline{the predicted number of cones sold increases by about 3.4.}

  \item Suppose the true population slope were $\beta = 0$. What would that mean
        about the relationship between temperature and ice cream sales?
        \ansline{No real linear relationship;}\\\ansline{temperature would not help predict sales; $b \ne 0$ only by chance.}

  \item Does $r = 0.87$ alone prove that $\beta \ne 0$? Explain.
        \ansline{No. $r$ is a statistic;}\\\ansline{it can be large by chance even if $\beta = 0$. A test is needed.}
\end{enumerate}
\end{practicebox}

\end{document}
